% ============================================================================
% APPENDIX (Mandatory for KDD)
% No page limit - reviewers not required to read
% ============================================================================

\appendix

\section{Reproducibility}
\label{app:reproducibility}

\subsection{Code and Data Availability}

All code, data, and pre-trained priors are available at:

\begin{center}
\url{https://github.com/redacted/banditgpt}
\end{center}

\paragraph{Repository Contents.}
\begin{itemize}[leftmargin=*]
    \item \texttt{banditgpt/} --- Core library implementation
    \item \texttt{banditgpt/data/priors/} --- Bundled expert priors (<1MB)
    \item \texttt{experiments/} --- Scripts to reproduce all figures and tables
    \item \texttt{kdd\_paper/} --- Camera-ready figures, tables, and LaTeX source
    \item \texttt{tests/} --- 127 unit and integration tests
\end{itemize}

\paragraph{Reproducing Results.}
{\small
\begin{verbatim}
# Install
pip install banditgpt

# Reproduce figures
python experiments/run_rq1.py  # Figure 1
python experiments/run_rq2.py  # Figure 3
python experiments/run_rq3.py  # Figure 4

# Reproduce latency benchmarks
python experiments/benchmark_latency.py
\end{verbatim}
}

\paragraph{Hardware and Software.}
\begin{itemize}[leftmargin=*]
    \item \textbf{Hardware}: Experiments run on a single CPU (Apple M2, 16GB RAM). No GPU required for routing inference.
    \item \textbf{Software}: Python 3.10+, NumPy, Sentence-Transformers, PyTorch (for embeddings only)
    \item \textbf{API Access}: Model responses generated via OpenRouter API (requires API key for regeneration; cached responses provided)
\end{itemize}

\paragraph{Random Seeds.}
All experiments use \texttt{seed=42} for reproducibility. Results may vary slightly ($<1\%$) due to API response variability.

% ============================================================================
% Ethical Use of Data
% ============================================================================

\section{Ethical Considerations}
\label{app:ethics}

\subsection{Data Sources}

This work uses only publicly available datasets and commercial APIs:

\begin{itemize}[leftmargin=*]
    \item \textbf{LMSYS Chatbot Arena}~\cite{zheng2023judging}: CC-BY-4.0 licensed. Used for archetype prompt clustering (497 prompts). No personally identifiable information (PII) is present in the sampled prompts.
    
    \item \textbf{OpenRouter API}: Commercial API for LLM inference. All model responses are generated via standard API calls with no special access or data collection agreements.
    
    \item \textbf{GPT-4o as Judge}: Used as an oracle grader for quality estimation. Prompts are standard evaluation queries with no sensitive content.
\end{itemize}

\subsection{Intended Use}

The BanditGPT router is designed for:
\begin{itemize}[leftmargin=*]
    \item Cost optimization in enterprise LLM deployments
    \item Quality-aware model selection for production systems
    \item Research on contextual bandit algorithms for LLM orchestration
\end{itemize}

The system does \textbf{not} perform content moderation, safety filtering, or policy enforcement. These concerns should be addressed by the underlying LLM providers or separate safety layers.

\subsection{Potential Risks}

\begin{itemize}[leftmargin=*]
    \item \textbf{Model Bias Propagation}: If the teacher oracle (GPT-4o) has biases, these may be encoded in the shipped priors. Mitigation: the online learning component allows the router to adapt to local feedback, potentially correcting biased priors.
    
    \item \textbf{Cost Optimization vs. Quality}: Aggressive cost optimization ($\lambdacost > 50$) may route queries to lower-quality models. Mitigation: we recommend the ``balanced'' profile for most deployments.
\end{itemize}

% ============================================================================
% Scalability Details
% ============================================================================

\section{Zero-Overhead Scalability: Technical Details}
\label{app:scalability}

This section elaborates on the scalability architecture that enables $O(1)$ model onboarding.

\subsection{Benchmarks as Metadata, Not Training Data}

A key architectural distinction is that we use benchmark scores only as \textbf{coarse safety filters}, not as training data for the router:

\begin{itemize}[leftmargin=*]
    \item \textbf{FrugalGPT/RouteLLM}: Require running calibration prompts through new models to train/calibrate the router. Adding Model X requires $\sim$500--1000 API calls to learn its behavior. This is $O(N)$ scaling.
    
    \item \textbf{BanditGPT}: The ``brain'' (prior covariances) learns what makes prompts \emph{difficult}---this knowledge is model-agnostic. A new model simply joins the pool with its public metadata. Time to onboard: seconds.
\end{itemize}

\subsection{The ``Warm Start'' vs.\ ``New Model'' Distinction}

\paragraph{Warm Start (One-Time, Offline).} We train the system on historical data to teach it what ``hard'' looks like. The 81-model $\times$ 497-archetype evaluation generates teacher supervision for prior fitting. This is a \textbf{fixed cost} amortized across all deployments.

\paragraph{New Model (Continuous, Online).} When a new model releases, it enters an ecosystem that is \emph{already intelligent}. The router knows: ``This prompt embedding pattern indicates hard math---I won't give it to the new model until it proves itself on easier prompts.'' The new model doesn't need its own ``warm start'' because the intelligence is in prompt understanding, not model enumeration.

\subsection{Safety Filters from Public Metadata}

We use publicly reported benchmark scores (e.g., from release announcements or Artificial Analysis) as \textbf{optional hard constraints}:

\begin{verbatim}
# Day 1: Register new model with metadata
router.add_model(
    model_id="new-model/v1",
    cost_per_1k=0.50,
    metadata={"math_score": 85}  # From release blog
)

# The bandit handles everything else
result = router.route(prompt, min_quality=80)
\end{verbatim}

The \texttt{min\_quality=80} filter ensures the new model is only eligible for prompts where the quality floor is below its claimed capability. The bandit then autonomously verifies whether the model \emph{actually} performs at that level via UCB exploration.

\subsection{Self-Correction via Online Learning}

If a model's claimed benchmarks are inflated, the bandit self-corrects:

\begin{enumerate}[leftmargin=*,nosep]
    \item Initial: Model X claims 90\% math accuracy; allocated exploration traffic.
    \item Reality: Actual performance is 70\% (reward signal from judge feedback).
    \item Update: Bandit reduces $\hat{Q}(q, X)$ via rank-one covariance updates.
    \item Convergence: Within $\sim$50--100 samples, traffic share reflects true capability.
\end{enumerate}

This ``trust but verify'' approach enables immediate deployment without blocking on expensive offline profiling, while the online learning ensures the system converges to accurate model assessments.

% ============================================================================
% Additional Tables (Overflow from Main Paper)
% ============================================================================

\section{Additional Experimental Results}
\label{app:tables}

\subsection{Full ROI Leaderboard}

Table~\ref{tab:roi_full} presents the complete ROI leaderboard for all evaluated models.

\begin{table}[h]
\centering
\footnotesize
\caption{\textbf{Full ROI Leaderboard.} All 81 models ranked by ROI Factor = $\|\boldsymbol{\theta}\| / \text{Cost}$ relative to GPT-4o baseline. Top 15 models shown.}
\label{tab:roi_full}
\resizebox{\columnwidth}{!}{%
\begin{tabular}{clccc}
\toprule
\textbf{Rank} & \textbf{Model} & \textbf{Cost/1k} & \textbf{$\|\theta\|$} & \textbf{ROI} \\
\midrule
1 & Nova-Micro-v1 & \$0.06 & 2.16 & 93.3× \\
2 & Nova-Lite-v1 & \$0.11 & 3.66 & 91.6× \\
3 & Llama-3.2-1B & \$0.05 & 1.48 & 73.6× \\
4 & Llama-3.2-3B & \$0.06 & 1.64 & 71.8× \\
5 & DeepSeek-R1-8B & \$0.07 & 1.30 & 50.5× \\
6 & Llama-3-8B & \$0.07 & 1.33 & 50.0× \\
7 & Gemini-Flash-Lite & \$0.13 & 1.58 & 31.7× \\
8 & Ministral-3B & \$0.10 & 1.15 & 30.2× \\
9 & Llama-3.2-11B-Vision & \$0.16 & 1.69 & 27.8× \\
10 & Llama-3.1-8B & \$0.10 & 0.98 & 25.9× \\
11 & GPT-OSS-20B & \$0.10 & 0.92 & 24.3× \\
12 & GPT-4.1-Nano & \$0.18 & 1.45 & 21.9× \\
13 & Ministral-8B & \$0.15 & 1.24 & 21.8× \\
14 & Gemini-2.0-Flash & \$0.18 & 1.41 & 21.1× \\
15 & Mistral-Small-24B & \$0.15 & 1.17 & 20.5× \\
\midrule
\textit{Ref} & \textit{GPT-4o} & \textit{\$4.38} & \textit{1.66} & \textit{1.0×} \\
\bottomrule
\end{tabular}%
}
\end{table}

\subsection{Detailed Latency Breakdown}

Table~\ref{tab:latency_full} provides detailed latency statistics across multiple runs.

\begin{table}[h]
\centering
\footnotesize
\caption{\textbf{Detailed Latency Benchmarks.} Router inference latency across 1,000 runs (81 models, 384 dims).}
\label{tab:latency_full}
\begin{tabular}{lrrrrr}
\toprule
\textbf{Prior Type} & \textbf{Mean} & \textbf{P50} & \textbf{P95} & \textbf{P99} & \textbf{Max} \\
\midrule
Expert (Disjoint) & 6.5ms & 6.4ms & 7.3ms & 8.9ms & 12.9ms \\
Shared Priors & 3.8ms & 3.8ms & 4.0ms & 4.1ms & 9.3ms \\
Cold Start & 3.2ms & 3.2ms & 3.5ms & 3.9ms & 7.1ms \\
\bottomrule
\end{tabular}
\end{table}

\subsection{Hyperparameter Sensitivity}

Table~\ref{tab:hyperparams} shows sensitivity to key hyperparameters.

\begin{table}[h]
\centering
\small
\caption{\textbf{Hyperparameter Sensitivity.} Effect of prior strength ($\lambda_{\text{boost}}$) and exploration rate ($\alpha$) on regret reduction vs.\ cold-start baseline.}
\label{tab:hyperparams}
\begin{tabular}{ccc}
\toprule
$\lambda_{\text{boost}}$ & $\alpha$ & \textbf{Regret Reduction} \\
\midrule
1.0 & 0.5 & 12.3\% \\
10.0 & 0.5 & 41.7\% \\
\textbf{50.0} & \textbf{0.5} & \textbf{63.6\%} \\
100.0 & 0.5 & 58.2\% \\
50.0 & 0.1 & 61.4\% \\
50.0 & 1.0 & 55.8\% \\
\bottomrule
\end{tabular}
\end{table}

% ============================================================================
% Algorithm Pseudocode
% ============================================================================

\section{Algorithm Details}
\label{app:algorithm}

\subsection{Routing Algorithm}

Algorithm~\ref{alg:routing} presents the complete routing procedure.

\begin{algorithm}[h]
\caption{Density-Based Warm-Start Routing}
\label{alg:routing}
\begin{algorithmic}[1]
\REQUIRE Query $q$, Model pool $\mathcal{M}$, Priors $\{(\mathbf{A}_m, \mathbf{b}_m)\}_{m \in \mathcal{M}}$
\REQUIRE Cost weights $\lambdacost, \lambdalat$, Exploration rate $\alpha$
\STATE $\mathbf{x} \leftarrow \text{Embed}(q)$ \COMMENT{384-dim sentence embedding}
\FOR{each model $m \in \mathcal{M}$}
    \STATE $\boldsymbol{\theta}_m \leftarrow \mathbf{A}_m^{-1} \mathbf{b}_m$ \COMMENT{Learned weights}
    \STATE $\mu_m \leftarrow \boldsymbol{\theta}_m^\top \mathbf{x}$ \COMMENT{Predicted quality}
    \STATE $\sigma_m \leftarrow \sqrt{\mathbf{x}^\top \mathbf{A}_m^{-1} \mathbf{x}}$ \COMMENT{Uncertainty}
    \STATE $\text{UCB}_m \leftarrow \mu_m + \alpha \cdot \sigma_m$ \COMMENT{Upper confidence bound}
    \STATE $U_m \leftarrow \text{UCB}_m - \lambdacost \cdot C_m - \lambdalat \cdot L_m$ \COMMENT{Utility}
\ENDFOR
\STATE $m^* \leftarrow \argmax_{m \in \mathcal{M}} U_m$
\RETURN $m^*$
\end{algorithmic}
\end{algorithm}

\subsection{Online Update}

After receiving feedback $(m, \mathbf{x}, r)$:

\begin{algorithmic}[1]
\STATE $\mathbf{A}_m \leftarrow \mathbf{A}_m + \mathbf{x}\mathbf{x}^\top$ \COMMENT{Reduce uncertainty}
\STATE $\mathbf{b}_m \leftarrow \mathbf{b}_m + r \cdot \mathbf{x}$ \COMMENT{Update reward direction}
\end{algorithmic}

This rank-one update takes $O(d^2)$ time, enabling real-time learning.

% ============================================================================
% SLA-Aware Routing (Moved from Main Paper)
% ============================================================================

\section{SLA-Aware Routing Framework}
\label{app:sla}

Rather than a rigid alternative to cascading, BanditGPT offers a \textbf{unified framework} controlled by a verification threshold ($\tau_{\text{verify}}$):

\begin{equation}
\text{Action} = \begin{cases} 
\text{Single-Shot} & \text{if } P(\text{Success}) > \tau_{\text{verify}} \\
\text{Cascade} & \text{if } P(\text{Success}) \le \tau_{\text{verify}}
\end{cases}
\end{equation}

\paragraph{Threshold Settings.}
\begin{itemize}[leftmargin=*]
    \item $\tau_{\text{verify}} = 0.0$ (Cost-Optimal): Never cascade (single-shot routing).
    \item $\tau_{\text{verify}} = 0.6$ (Balanced): Cascade when uncertain.
    \item $\tau_{\text{verify}} = 1.0$ (Max Assurance): Always cascade (equivalent to always verifying).
\end{itemize}

\paragraph{Hard Budget Constraints.} BanditGPT supports hard constraints via a constrained action space:
\begin{equation}
\mathcal{A}_{\text{feasible}} = \{m \in \mathcal{M} : \text{cost}(m) \le C_{\max}, \text{latency}(m) \le L_{\max}\}
\end{equation}

When constrained to \$0.50/1k: GPT-4o and DeepSeek fail the constraint entirely, while BanditGPT maintains 85\% accuracy by auto-routing to cheap specialists.

\begin{figure}[h]
    \centering
    \includegraphics[width=\linewidth]{figures/figure7_sla_tunability.pdf}
    \caption{\textbf{Dialing the Pareto Frontier.} Users dynamically trade cost for accuracy via $\tau_{\text{verify}}$.}
    \label{fig:sla_tunability_app}
\end{figure}

\begin{figure}[h]
    \centering
    \includegraphics[width=\linewidth]{figures/figure8_finops_constraints.pdf}
    \caption{\textbf{Resilience Under Hard Budget Constraints.} Static policies fail; BanditGPT adapts.}
    \label{fig:finops_app}
\end{figure}
